\section{Ampliación a sistemas modernos: leer la arquitectura de un agent a partir de Cicero}

Esta sección es una interpretación de ingeniería hecha en 2026 y no forma parte de los hechos de la clase de 2023. No presenta productos posteriores como argumentos del ponente. Traduce las interfaces que Cicero ya mostró a un lenguaje más general de diseño de agents.

\subsection{Separar «qué decir» de «qué hacer»}

Los agents modernos suelen mantener a la vez un private plan, tool actions y una user-facing response. La experiencia de Cicero muestra que entregar el plan interno completo a un modelo de lenguaje para que lo exprese libremente produce filtraciones de información, contradicciones y problemas por falta de capacidad de auditoría. Es mejor formar primero un action intent estructurado y dejar que una communication layer restringida decida qué información se comunica al exterior.

\begin{knowledgebox}{Interfaz transferible de cinco capas}
\begin{enumerate}
  \item \textbf{State estimator}: convierte el entorno, el historial y la retroalimentación en la belief actual;
  \item \textbf{Policy prior}: propone candidatos comunes, de bajo costo y compatibles con las convenciones humanas;
  \item \textbf{Planner / search}: asigna inference compute adicional según el valor de la tarea;
  \item \textbf{Intent interface}: comprime las decisiones internas en una estructura verificable;
  \item \textbf{Communication evaluator}: predice cómo cambiará la salida el entorno y el comportamiento de los demás.
\end{enumerate}
\end{knowledgebox}

\subsection{Estructura común de la prompt injection y los mensajes estratégicos}

Los mensajes de los oponentes en Diplomacy no son instrucciones confiables del sistema. Son una observation del entorno que persigue un objetivo propio. Esto se parece al límite de seguridad que enfrenta un agent que usa herramientas ante páginas web, correos o documentos de terceros que no son confiables: el contenido puede actualizar la belief, pero no debe obtener el control de forma automática. El enfoque de Cicero, «predecir el comportamiento del otro y luego juzgarlo con su propio value», ofrece un paradigma más robusto que el instruction-following directo.

\begin{warningbox}{Los límites de la analogía}
Diplomacy tiene reglas cerradas, acciones que pueden simularse y una puntuación explícita. Los agents de mundo abierto suelen carecer de un transition model y de una reward confiables. Por ello no basta con copiar piKL o el value filter. Lo que debe heredarse es su idea de auditoría: distinguir los datos de las instrucciones, registrar el intent de forma explícita y hacer que las salidas de alto riesgo pasen por una evaluación independiente.
\end{warningbox}

\subsection{Gestión de costos: no toda solicitud amerita search}

El cómputo en tiempo de inferencia debe tener una política de presupuesto. Un sistema práctico necesita estimar el valor del problema, la dificultad de verificación, el costo de un fallo y la latency restante. Con eso decide el número de candidatos, la profundidad de búsqueda, las tool calls y la intensidad del evaluator. Sin una condición de parada, el agent convierte «pensar más» en un bucle infinito. Si toda solicitud se resuelve con una sola greedy decode fija, se renuncia a la asignación dinámica de recursos, que es lo más valioso de search.

\begin{importantbox}{El problema del presupuesto en ingeniería}
Cada rollout o llamada al evaluator adicional debe poder responder cuatro preguntas: qué tipo de error reduce, si existe una señal de verificación independiente, cuándo el beneficio marginal se vuelve negativo y si la correlación entre errores puede invalidar la votación por mayoría. La inference compute solo es capacidad del sistema, y no un aumento de la factura, cuando entra en un proceso de decisión medible.
\end{importantbox}

\subsection{Resumen de la sección}

El legado más valioso de Cicero para los agents modernos no es un modelo concreto, sino los contratos entre módulos. El lenguaje que no es confiable entra primero en la belief y no controla las acciones de forma directa. El planning produce un intent estructurado. El lenguaje dirigido al exterior se somete, antes de enviarse, a una evaluación de sus consecuencias en el comportamiento. El presupuesto de cómputo varía según el valor de la tarea.
